\documentclass{article}

%%% Package List

\usepackage{datetime} % Dates.
\usepackage{fancyhdr} % Headers.
\usepackage{lastpage} % Page References.

\usepackage{graphicx} % Inserting Images.
\usepackage{float} % Postitioning Images.

\usepackage{dirtytalk} % Speech & Quote Marks.
\usepackage{hyperref} % Hyperlinks and References.
\usepackage{xcolor} % Coloured Text.
\usepackage{enumitem} % Letter Bullets. 

\usepackage{amssymb} % Maths Symbols.
\usepackage{amsmath} % Maths Tools.

\usepackage{algpseudocode} % Pseudocode.
\usepackage{algorithm} % Algorithms.
\usepackage{minted} % Code Snippets.

\usepackage[margin=1in]{geometry} % Reduces the Page Margins.

%%%



%%% Custom Formats and Commands

\newcommand{\templateversion}{1.5.1} % Do not edit this variable.

\hypersetup{colorlinks=true, urlcolor=black, linkcolor=black}


% Date Formats.
\newdateformat{monthyeardate}
{%
    \monthname[\THEMONTH] \THEYEAR
}

\newdateformat{yeardate}
{%
    \THEYEAR
}


% Colouring and Spacing.
\newcommand{\colour}[2] {\color{#1}#2\color{black}}
\newcommand{\separate}[1] {\vspace{#1 pt}\noindent}


% Maths Number Sets.
\newcommand{\R}{\mathbb{R}}
\newcommand{\N}{\mathbb{N}}
\newcommand{\C}{\mathbb{C}}
\newcommand{\Q}{\mathbb{Q}}
\newcommand{\Z}{\mathbb{Z}}
\newcommand{\I}{\mathbb{I}}


% Custom Formats.
\newcommand{\definition}[2]{\textbf{Definition #1}: #2\label{def:#1}}
\newcommand{\example}[2]{\textbf{Example #1}: #2\label{ex:#1}}
\newcommand{\conjecture}[2]{\textbf{Conjecture #1}: #2\label{cnj:#1}}
\newcommand{\numberedentry}[3]{\textbf{#2 #1}: #3\label{entry:#1}}
\newcommand{\proof}[2]{\textit{Proof}: #2 \begin{flushright}\(\square\)\end{flushright}\label{prf:#1}}
\newcommand{\theorem}[3]{\textbf{Theorem #1}: #2 \par\separate{3}\proof{thm:#1}{#3}\label{thm:#1}}
\newcommand{\lemma}[3]{\textbf{Lemma #1}: #2 \par\separate{3}\proof{lem:#1}{#3}\label{lem:#1}}


% Join Symbols.
\newcommand{\ojoin}{\setbox0=\hbox{$\bowtie$}%
    \rule[-.02ex]{.25em}{.4pt}\llap{\rule[\ht0]{.25em}{.4pt}}}
\newcommand{\insertinnerjoin}{\mathbin{\bowtie\mkern-5.8mu}\;}
\newcommand{\insertleftouterjoin}{\mathbin{\ojoin\mkern-5.8mu\bowtie}}
\newcommand{\insertrightouterjoin}{\mathbin{\bowtie\mkern-5.8mu\ojoin}}
\newcommand{\insertfullouterjoin}{\mathbin{\ojoin\mkern-5.8mu\bowtie\mkern-5.8mu\ojoin}}


% Advanced Formats.
\newcommand{\insertimage}[3]
{
    \begin{figure}[H]
        \centering
        \includegraphics[width=#3\linewidth]{#1}
        \caption{#2}
        \label{fig:#2}
    \end{figure}
}

\newcommand{\insertalgorithm}[4]
{
    \begin{algorithm}[H]
        \caption{#1}
        \begin{algorithmic}
            \State \textbf{Input}: #2
            \State \textbf{Output}: #3
            \separate{7}#4
        \end{algorithmic}
        \label{alg:#1}
    \end{algorithm}
}

\newcommand{\inserttable}[4]
{
    \begin{table}[H]
        \centering
        \begin{tabular}{|#1|}
            \hline #3 \\ \hline #4 \\ \hline
        \end{tabular}
        \caption{#2}
        \label{tab:#2}
    \end{table}
}

%%%



%%% Document Setup

% Title Page.
\title{\thetitle\\[0.5em]\large{\thesubtitle}}
\author{\theauthor}
\date{\thedate} 


% Custom Page Header.
\fancyhead{}
\fancyhead[R]{\thetitle\text{ - }\thesubtitle}

% Custom Page Footer.
\fancyfoot{}
\fancyfoot[L]{Written by \theauthor \\ 
\textit{\LaTeX \,Template v.\templateversion: Copyright \date{\yeardate\today}, Matthew Emerson}} % Do not edit.
\fancyfoot[R]{Page \thepage \, of\, \pageref{LastPage}}

%%%



%%% Page Variables

\newcommand{\thetitle}{Computer Systems} % Insert the document title here.
\newcommand{\thesubtitle}{COMP1071 - Durham University} % Insert the document subtitle here.
\newcommand{\theauthor}{Matthew Emerson} % Insert the name(s) of the author(s) here.
\newcommand{\thedate}{\monthyeardate\today} % Replace this if you want a static date.

%%%



%%% Setting Up Pages.

\pagestyle{empty}\begin{document}
\maketitle\thispagestyle{empty} % Title Page.
\newpage\tableofcontents % Contents Page.

\clearpage\setcounter{page}{1}\newpage % Excludes the title and contents page from page count.
\pagestyle{fancy}

%%%



%%% Document Content

\part{Digital Electronics}\label{part:1.de}

\section{Circuits}

\subsection{Transistors}

\definition{1.1.1}{\underline{Transistors} are electronically controlled switches that turn on or off when voltage or current is applied to a control terminal.}

\separate{7}\definition{1.1.2}{A \underline{nMOS transistor} is a sandwich of layers of conducting and insulating material. The gate is a conducting layer, on top of an insulating layer, on top of the substrate which is a semi-conductor. When its gate \(g\) is at a high voltage, current is allowed to flow from source to drain, and this is prevented when the gate is at a low voltage.}

\insertimage{Images/Transistors - nMOSTransistor.png}{An nMOS transistor}{0.5}

\noindent The junction between p-type and n-type silicon is a diode. This means the current can only flow from p-type to n-type as they act as anodes and cathodes respectively.

\separate{7}\definition{1.1.3}{A \underline{pMOS transistor} is the opposite of a nMOS transistor - it allows current to flow when the gate is at low voltages, and prevents current flow at high voltages.}

\separate{7}\definition{1.1.4}{The \underline{static discipline} states that if inputs to transistors meet valid thresholds, then the system must guarantee outputs will meet valid output thresholds and will not fall into the forbidden zone (outside of the accepted range for high/low voltages).}

\separate{7}\definition{1.1.5}{\underline{Moore's Law} is usually quoted as \say{transistor density doubles every two years}, or \say{computer processing power doubles every eighteen months}.}

\subsection{Circuit Design}

\separate{7}\definition{1.2.3}{\underline{Abtraction} is used when designing circuits to simplify problems by hiding unimportant details, highlighting the actual logic problem to be solved.}

\separate{7}\definition{1.2.4}{Another design principle used is \underline{discipline}. Here, design choices are deliberately restricted to make circuits easier to model, design, and combine.}

\separate{7}\definition{1.2.5}{When designing complex systems, engineers use \underline{the three-y's}: \underline{hierarchy}, dividing a system into clear modules and submodules; \underline{modularity}, using well-defined functions and interfaces for modules; and \underline{regularity}, encouraging uniformity so modules can be swapped or reused.}

\newpage\subsection{Combinatorial Circuits}

\definition{1.3.1}{A \underline{combinatorial circuit} is a circuit where the output depends purely on the current input.}

\separate{7}\definition{1.3.2}{An \underline{adder circuit} takes the contents of two registers and adds them together.}
\begin{itemize}
    \item \definition{1.3.3}{A \underline{half-adder} circuit takes two one-bit binary inputs and adds them, returning a sum bit, and a carry bit.}
    \item \definition{1.3.4}{A \underline{full-adder} circuit takes three one-bit binary inputs (one of which is a carry bit) and adds them, returning a sum bit, and a carry bit.}
\end{itemize}

\separate{7}\definition{1.3.5}{A \underline{decoder} circuit takes an n-bit binary number and decodes it into \(2^n\) data lines.}

\separate{7}\definition{1.3.6}{A \underline{multiplexor} (mux) circuit uses a binary number to select one of a set of inputs to funnel to its output.}

\separate{7}\definition{1.3.7}{A \underline{tristate} takes in two bits, an input bit and an enable bit. If the enable bit is a one, the input bit is outputted, if not then the output is left floating.}

\subsection{Sequential Circuits}

\definition{1.4.1}{A \underline{sequential circuit} is a circuit which changes its output depending on its input history, alongside the current input.}

\separate{7}\definition{1.4.2}{A circuit which has two stable output states for a given input is known as a \underline{bistable} circuit.}

\separate{7}\definition{1.4.3}{An \underline{sr-latch} takes in two inputs, set \(S\) and reset \(R\). If set has a pulse of one, then the output remains at one, even after the pulse ends. If reset has a pulse of one, then the output is reset back to zero. This is a bistable circuit.}

\separate{7}\definition{1.4.4}{A \underline{d-latch} has two inputs, a data input \(D\) and a clock input \(CLK\). The latch only changes to match the input at \(D\) when the \(CLK\) signal is high.}

\separate{7}\definition{1.4.5}{A \underline{d-type flip-flop} works similarly to the d-latch, but only allows the stored output to change on the rising edge of the clock signal.}

\section{Introduction to Logic}

\subsection{Functionally Complete Sets}

\definition{2.0.1}{A \underline{functionally complete set} is a set of boolean operators which can be used to express all possible truth tables by combining members of the set into a boolean expression. For example,
\(\{AND,\,OR,\,NOT\},\;\{NOR\},\text{ and}\;\{NAND\}\) are all functionally complete sets.}

\newpage\part{Machine Architecture}\label{part:2.ma}

\section{MIPS}

\subsection{Introduction to MIPS}

\definition{3.1.1}{A \underline{MIPS (Microprocessor without Interlocked Pipeline Stages)} is 32-bit RISC processor used in embedded systems and consumer electronics. It uses 32-bit words, has around 110 instructions, and 32 general purpose registers.}

\separate{7}\definition{3.1.2}{\underline{Amdahl's Law} states that the performance improvement to be gained from using some faster mode of execution is limited by a fraction of the time the faster mode can be used.}

\subsection{The MIPS Instruction Set}

The instruction set uses five addressing modes, these are:
\begin{itemize}
    \item \definition{3.2.1}{When all addresses in the instruction are registers, this is \underline{register-only addressing}.}
    \item \definition{3.2.2}{When a value is used directly in an instruction, this is \underline{immediate addressing}.}
    \item \definition{3.2.3}{\underline{Base addressing} is when the address of a memory operand is given by a base address and a signed immediate.}
    \item \definition{3.2.4}{\underline{PC-relative addressing} takes the current position in memory and jumps a given amount to a new address.}
    \item \definition{3.2.5}{\underline{Pseudo-direct addressing} is used by MIPS in place of direct addressing as it does not support 32-bit addresses. As a result, it calculates it from 26 as the two least significant bits are always zeros as all words are multiples of four, and the four most significant bits can also be left as zero.}
\end{itemize}

\separate{7}The MIPS instruction set can be split into three instruction types, these are:
\begin{itemize}
    \item \definition{3.2.6}{\underline{R-type instructions} contain a six-bit opcode, three five-bit registers (two source and one destination), a five-bit shift value, and a six-bit code to identify the function.}
    \item \definition{3.2.7}{\underline{I-type instructions} contain a six-bit opcode, two five-bit registers (one source, one destination), and a sixteen-bit immediate.}
    \item \definition{3.2.8}{\underline{J-type instructions} contain a six-bit opcode and a twenty-six bit memory address. The compiler combines the two to form a new thirty-two bit address.}
\end{itemize}

\subsection{Memory Maps}

MIPS uses 32-bit memory addresses across \(2^32\) bits (or 4 GB), ranging from 0x0 to 0xFFFFFFFC. This address space is split into the four segments:
\begin{itemize}
    \item \definition{3.3.1}{The \underline{text segment} stores the machine language program and can store 256 MB.}
    \item \definition{3.3.2}{The \underline{global data segment} stores global variables for program functions to access. These can be accessed using the pointer \(\$gp\). This segment can store 64 KB of data.}
    \item \definition{3.3.3}{The \underline{dynamic data segment} stores data which is dynamically allocated/deallocated through the execution of the program and can store up to 2 GB of data using a stack and a heap.}
    \item \definition{3.3.4}{The \underline{reserved segment} is used by the operating system and cannot be directly accessed by the program.}
\end{itemize} 

\section{The ATmega328P Microcontroller}

\subsection{Main Technical Characteristics}

\definition{4.1.1}{A \underline{microcontroller} is an integrated device consisting of a microprocessor, memory, and an input/output. They are used in embedded systems and store their program in flash memory.}

\separate{7}The microcontroller has 32 general purpose 8-bit registers that are copied onto SRAM memory. Registers 26-31 are used in addressing, and as such should be treated as three 16-bit registers, rather than six 8-bit registers.

\separate{7}\definition{4.1.2}{The \underline{status register} is an 8-bit register separate from the other 32. Its eight bits are used as flags.}

\separate{7}The microcontroller's pins are separated into 3 groups, B (digital), C (analogue), and D (digital). 

\separate{7}\definition{4.1.3}{The direction of each pin of a port is controlled by a \underline{data direction registers}. We say a bit is set when it is 1 and the corresponding pin is output, and that a bit is cleared when it is 0 and the pin is input.}

\section{Programming in AVR Assembly}

\subsection{Directives}

\definition{5.1.1}{The assembly code contains statements called \underline{directives}. These do not correspond to assembly instructions, but the assembler will take them into account when compiling. In AVR assembly, a directive is preceded by a dot. Some examples include:}
\begin{itemize}
    \item \(.equ\) assigns a value to a label which can be used in later expressions. This is a constant. (\(.equ\;pi, 3\)).
    \item \(.global\;start\) and \(.global\;btnLED\) declare the functions \(start\) and \(btnLED\) as global.
\end{itemize}

\newpage\part{Databases}\label{part:3.db}

\section{Database Design}

\subsection{The ANSI-SPARC Architecture}

The architecture, set out by ANSI (the American National Standards Institute) and SPARC(the Standards, Planning, and Requirements Committee), consists of three levels used to realise data abstraction allowing different users to access the same data in different views without affecting each other. These levels are:
\begin{itemize}
    \item \definition{6.1.1}{The \underline{external layer} is the user's view of the database. It can be adapted by the user to fit their specific needs without affecting others. This layer also encorporates derived, calculated, or combined fields.}
    \item \definition{6.1.2}{The \underline{conceptual layer} is the community view of the database. It describes what data is stored in the database alongside the logical structure of the entire database.}
    \item \definition{6.1.3}{The \underline{internal layer} is the physical representation of the database which describes how the data is stored to achieve optimal performance and storage utilisation.}
\end{itemize}

\separate{7}\definition{6.1.4}{A \underline{database schema} is a complete description of the database as a result of the database design process. Each database should contain multiple external schemata, a conceptual schema, and an internal schema.}

\separate{7}\definition{6.1.5}{A \underline{database instance} is the concrete data contained within the database at any particular time.}

\separate{7}\definition{6.1.6}{A database is considered to have \underline{logical data independence} if the external schemata remain the same if the logical structure is changed on a conceptual level.}

\separate{7}\definition{6.1.7}{A database is considered to have \underline{physical data independence} if the conceptual schema stays the same when the internal schema is changed.}

\subsection{Desigining a Database}

\definition{6.2.1}{When creating a database, the first thing to do is to create a \underline{conceptual design}. This is a high-level model of the data using entity-relationship modelling, identifying the appropriate entities, attributes, relationships, and constraints. This should be independent of any physical constraints.}

\separate{7}\definition{6.2.2}{The next step is to create a \underline{logical design} which will be used to construct the relational data model using the conceptual design as a guide to transform entities and relationships into tables, using normalisation techniques to eliminate redundancies.}

\separate{7}\definition{6.2.3}{The \underline{physical design} should be created last, and will describe how the logical design will be implemented, defining storage, access, and security details.}

\newpage\section{Modelling Databases}

\subsection{Entity-Relationship Modelling}

\definition{7.1.1}{An \underline{entity} should be explicitly represented, denoted by a labelled rectangle and split into types and instances.}

\separate{7}\definition{7.1.2}{A \underline{relationship} is a named association between two (or more) entity types which is represented by a labelled line with a triangle indicating reading direction.}

\separate{7}\definition{7.1.3}{\underline{Constraints} are commonly represented on an entity-relationship model by a range \((min...max)\). This can also include the cardinality (one-to-one, many-to-one, etc.), and participation (optional or mandatory).}

\separate{7}\definition{7.1.4}{The properties of an entity type are known as its \underline{attributes}. In an entity-relationship model they are placed below an entities title with the primary key underlined.}

\subsection{Relational Data Models}

\definition{7.2.1}{A \underline{relational data model} is based on the concept of mathematical relations and is the most widely used data structure for designing databases. It models a database by storing the relations between data in tables.}

\separate{7}When creating a relational data model, the design must meet the following integrity constraints:
\begin{itemize}
    \item \definition{7.2.2}{\underline{Domain constraints} for attributes limit the range of values permitted.}
    \item \definition{7.2.3}{\underline{Entity integrity} ensures that the primary key is not null, guaranteeing that each entity has a unique identifier and that foreign keys form a valid reference.}
    \item \definition{7.2.4}{\underline{Referential integrity} ensures that all foreign keys match some primary key or is null. This ensures that all references are valid.}
\end{itemize}

\subsection{Functional Dependencies}

\definition{7.3.1}{A relation's \underline{functional dependencies} describe the relationship among its attributes.}

\separate{7}\definition{7.3.2}{If \(A\) and \(B\) are two sets of attributes in the same relation, where \(A\) is associated with exactly one value of \(B\), then \(A\) \underline{functionally determines} \(B\) and \(B\) is \underline{functionally dependent} on \(A\). \(A\) is called the \underline{determinant} of \(B\). Functional dependencies can take one of three forms:}
\begin{itemize}
    \item \definition{7.3.3}{A \underline{full functional dependency} exists when \(B\) is functionally dependent on \(A\), but not on any proper subset of \(A\).}
    \item \definition{7.3.4}{A \underline{partial functional dependency} exists when \(B\) is functionally dependent on at least one proper subset of \(A\).}
    \item \definition{7.3.5}{A \underline{transitive functional dependency} exists between \(A\) and \(C\) if \(A \rightarrow B\) and \(B \rightarrow C\).}
\end{itemize}

\separate{7}\definition{7.3.6}{Given a set of functional dependencies \(F\), its closure \(F^+\) are all the functional dependencies which are implied by the dependencies in \(F\). These can be calculated using \hyperref[def:7.3.7]{Armstrong's Axioms}.}

\newpage\noindent\definition{7.3.7}{\underline{Armstrong's Axioms} are: 
\begin{itemize}
    \item \definition{7.3.8}{\underline{Reflexivity}: If \(B \subseteq A\), then \(A \rightarrow B\).}
    \item \definition{7.3.9}{\underline{Augmentation}: If \(A \rightarrow B\), then \((A,C) \rightarrow (B,C)\).}
    \item \definition{7.3.10}{\underline{Transivity}: If \(A \rightarrow B\) and \(B \rightarrow C\), then \(A \rightarrow C\).}
\end{itemize}} 
These rules are complete and sound, and can be used to derive another three rules:
\begin{itemize}
    \item \definition{7.3.11}{\underline{Decomposition}: If \(A \rightarrow (B,C)\), then \(A \rightarrow B\) and \(A \rightarrow C\).}
    \item \definition{7.3.12}{\underline{Union}: If \(A \rightarrow B\) and \(A \rightarrow C\), then \(A \rightarrow (B,C)\).}
    \item \definition{7.3.13}{\underline{Composition}: If \(A \rightarrow B\) and \(C \rightarrow D\), then \((A,C) \rightarrow (B,D)\).}
\end{itemize}

\separate{7}The closure \(F^+\) of \(F\) can be computed by taking of \(F\)'s dependencies, repeatedly applying Armstrong's Axioms, and adding any new dependencies discovered in this way.

\subsection{Normalisation}

\definition{7.4.1}{A table is in \underline{first normal form} if it has no repeating groups (every cell has one value).}

\separate{7}\definition{7.4.2}{A table is in \underline{second normal form} if it is in first normal form and has no partial dependencies (every non-key attribute is dependent on the whole primary key).}

\separate{7}\definition{7.4.3}{A table is in \underline{third normal form} if it is in second normal form and has no transitive dependencies.}

\section{Relational Calculus and Algebra}

\subsection{Database Languages}

\definition{8.1.1}{\underline{Structured query language} is the most common database language. It follows the ISO/ANSI standard using simple syntax with its two components, DDL and DML:}
\begin{itemize}
    \item \definition{8.1.2}{\underline{Data definition language} is used to define the schema for each relation, the domain for each respective attribute, and to specify the relation's integrity constraints.}
    \item \definition{8.1.3}{\underline{Data manipulation language} is used to insert, update, delete, and retrieve data from the database.}
\end{itemize}

\separate{7}Database languages can be separated into one of two types:
\begin{itemize}
    \item \definition{8.1.4}{A \underline{declarative language} specifies which data needs to be retrieved in a query (relational calculus).}
    \item \definition{8.1.5}{A \underline{procedural language} specifies how to recieve the requested data for the query (relational algebra).}
\end{itemize}

\newpage\subsection{Relational Calculus}

\definition{8.2.1}{\underline{Relational calculus} uses set theoretic expressions to define new sets based on one or more existing ones.}

\separate{7}\definition{8.2.2}{A \underline{predicate} is a function that returns either true or false when given a set of arguments.}

\separate{7}\definition{8.2.3}{A \underline{proposition} is the expression obtained when replacing values for the arguments of a predicate.}

\separate{7}For example, for the predicates \(P(x)\) and \(Q(x)\), the set such that both \(P(x)\) and \(Q(x)\) are true is 
\begin{equation}
    \{x|P(x) \land Q(x)\}
\end{equation}

\subsection{Relational Algebra}

\definition{8.3.1}{\underline{Relational algebra} uses operators to construct new sets from one or more existing ones.}

\separate{7}Relational algebra uses a unique set of operations and symbols. These include:
\begin{itemize}
    \item \definition{8.3.2}{The \underline{selection} operation, \(\sigma_c(R)\), is a unary operation which selects the subset of tuples from a relation \(R\) that meet the specified condition \(c\).}
    \item \definition{8.3.3}{The \underline{projection} operation, \(\Pi_{c_1,c_2}(R)\), is a unary operation which selects a specified subset of attributes \(c_1, c_2\) from a relation \(R\).}
    \item \definition{8.3.4}{The \underline{cartesian product} operation, \(A \times B\), is a binary operaiton which generates all possible combinations of tuples from \(A\) and \(B\) by concatenating them.}
    \item \definition{8.3.5}{The \underline{rename} operation, \(\rho_{Y(A,B,C)}(X)\), is a unary operation which takes the relation \(X\) and renames it to \(Y\), and naming its corresponding attributes \(A, B, C\).}
    \item \definition{8.3.6}{The \underline{division} operation, \(A \div B\), is a binary operation that only keeps tuples from \(A\) that exist in all possible combinations with \(B\) where \(B\) contains a subset of \(A\).}
    \item There are also join operations, corresponding to their SQL eqivalents: inner \(\insertinnerjoin\), left outer \(\insertleftouterjoin\), right outer \(\insertrightouterjoin\), and full outer \(\insertfullouterjoin\).
    \item Some other basic operations include \(A - B\), \(A \cup B\), and \(A \cap B\).
\end{itemize}

\newpage\part{Operating Systems}\label{part:4.os}

\section{Introduction to Operating Systems}

\subsection{Basic Definitions}

\definition{9.1.1}{An \underline{operating system} is a program that acts as an intermediary between a user and the hardware. Its goal is to execute programs, make solving problems easier, make the computer system convenient to use, and to use the resources of the system fairly and efficiently. The operating system is:}
\begin{itemize}
    \item \definition{9.1.2}{A \underline{resource allocator}. It is responsible for the management of the computer system's resources.}
    \item \definition{9.1.3}{A \underline{control program}. It controls the execution of user programs and operation of the input/output devices to prevent errors and improper use of the computer.}
    \item \definition{9.1.4}{A \underline{kernel}. It is the one program that runs all of the time. It consists of: the privileged instruction set, an interrupt mechanism, memory protection, and a real-time clock. It also has:}
    \begin{itemize}
        \item \definition{9.1.5}{A \underline{first-level interrupt handler} which determines the source of an interrupt and initiates servicing the interrupt.}
        \item \definition{9.1.6}{The \underline{dispatcher} which assigns processing resources for processes and is initiated when a current process cannot continue, or if the CPU may be used better elsewhere.}
        \item Intra operating system communications.
    \end{itemize}
\end{itemize}

\subsection{Operating System Design}

\separate{7}\definition{9.2.1}{The operating system generally runs in \underline{dual-mode} operation. This allows the operating system to protect itself and other system components. It differentiates which mode it is currently in using the \underline{mode bit}.}
\begin{itemize}
    \item \definition{9.2.2}{When the computer is in \underline{user mode} it is limited as to which instructions it can execute, with some requiring switching modes to,}
    \item \definition{9.2.3}{\underline{Kernel mode} (also known as system or privileged mode). In this mode, the system is given access to a set of \underline{privileged instructions} which can perform tasks such as managing interrupts, performing inputs/outputs, and halting a process.}
\end{itemize}

\separate{7}The operating system offers a set of services, some of which are:
\begin{itemize}
    \item Providing an environment for the execution of programs and services.
    \item Providing a user interface, which can either be on the command line (CLI), graphical (GUI), or in Batch.
    \item Controlling the operation of input/output devices.
    \item Manipulating the file system.
    \item Organising communications either between processes on the same computer, or between computers over a network.
    \item Providing error detection, taking appropriate action to ensure correct and consistent computing.
    \item Allocating resources, particularly when multiple jobs are running concurrently.
    \item Providing protection and security by ensuring that all access to system resources are controlled and that all users are authenticated.
\end{itemize}

\subsection{Processes and Threads}

\definition{9.3.1}{A program in execution is known as a \underline{process}. A process includes code, its current activity, a stack, a data section, and a heap.}

\separate{7}\definition{9.3.2}{A process's information is stored in a \underline{process control block}. These PCBs are then arranged together into a \underline{process table}. Each PCB stores a process's unique identifier, state, memory useage, CPU scheduling information, and its corresponding CPU registers.}

\separate{7}\definition{9.3.3}{As a process is executed, it changes state. Here, it can take one of five states: new, when the process is being created; running, when instructions are being executed; waiting, when the process is waiting for some event to occur; ready, when the process is ready to be dispatched to the CPU; and terminated, when the process has finished executing.}

\separate{7}\definition{9.3.4}{A process can, as a \underline{parent process}, create a number of \underline{child processes}, which in turn can create other processes, forming a \underline{tree of processes}. A parent process may share some of its resources with its children, and they can either execute concurrently or sequentially.}

\separate{7}\definition{9.3.5}{The \underline{process scheduler} selects which of the currently available processes will be executed next by the CPU. It maintains the following schedulling queues:}
\begin{itemize}
    \item \definition{9.3.6}{The \underline{job queue} is the set of all processes in the system.}
    \item \definition{9.3.7}{The \underline{ready queue} is the set of all processes residing in main memory, ready and waiting to execute.}
    \item \definition{9.3.8}{The \underline{device queues} are sets of processes waiting for I/O devices.}
\end{itemize}

\separate{7}\definition{9.3.9}{A \underline{thread} is a unit of execution which can be traced and that belongs to a process. It can either be a user or kernel thread.}

\separate{7}The thread model of processing uses local variables within threads, with global variables being available to all threads. The model is frequently used as: it requires less time to create or terminate threads vs. processes; it may allow for execution to continue, even if part of a process is blocked; threads can share processing resources; creating threads is cheaper than creating processes or context switching; and processes can take advantage of multiprocessor architectures, increasing scalability.

\separate{7}\definition{9.3.10}{\underline{Amdah's Law} states that performance gains from adding additional cores with serial and parallel components follows the following equation:
\begin{equation}
    speedup \leq \frac{1}{S+\frac{(1-S)}{N}}
\end{equation}
where \(S\) is the serial portion and \(N\) is the number of processing cores.}

\separate{3}For example, if an application is \(75\%\) parallel, moving from one to two cores will result in a speedup of 1.6 times. As \(N\) approaches infinity, the speedup approaches \(\frac{1}{S}\).

\separate{7}Multithreading can be accomplished by using one of three models:
\begin{itemize}
    \item \definition{9.3.11}{The \underline{many-to-one model} maps many user-level threads to a single kernel thread.}
    \item \definition{9.3.12}{The \underline{one-to-one model} maps each user-level thread to a kernel thread, allowing for greater concurrency.}
    \item \definition{9.3.13}{The \underline{many-to-many model} allows many user-level threads to be mapped to many kernel threads.}
\end{itemize}

\section{Process, Memory, and Storage Management}

\subsection{CPU Scheduling}

\definition{10.1.1}{The aim of \underline{multiprogramming} is to maximise CPU utilisation by keeping multiple processes in memory concurrently, meaning the CPU can always be executing a ready process whilst another is waiting.}

\separate{7}\definition{10.1.2}{The \underline{long-term (job) scheduler} selects which processes should be brought into the ready queue.}

\separate{7}\definition{10.1.3}{The \underline{medium-term scheduler} is a part of swapping and is in-charge of handling the swapped out processes.}

\separate{7}\definition{10.1.4}{The \underline{short-term (CPU) scheduler} selects the process to be executed next and allocates it to the CPU.}

\separate{7}\definition{10.1.5}{The \underline{dispatcher} is the module that gives control of the CPU to the process selected by the CPU scheduler. It is in charge of context switching, switching to user mode, and jumping to the proper location in the user program to restart that program.}

\separate{7}\definition{10.1.6}{The \underline{dispatch latency} is the time it takes for the dispatcher to stop one process and start another running.}

\separate{7}When scheduling processes, the criteria considered are:
\begin{itemize}
    \item CPU utilisation.
    \item \definition{10.1.7}{\underline{Throughput}, the number of processes that complete their execution per unit of time.}
    \item \definition{10.1.8}{\underline{Turnaround time}, the amount of time to execute a particular process (waiting, execute, and I/O).}
    \item \definition{10.1.9}{\underline{Waiting time}, the total amount of time a process has been waiting in the ready queue.}
    \item \definition{10.1.10}{\underline{Response time}, the amount of time it takes from when a request was submitted until the first response.}
\end{itemize}
This can be achieved using one of the following algorithms: first come, first served; shortest job first; shortest remaining time first; priority; round-robin; or multilevel/feedback queue.

\subsection{Memory Management}

\definition{10.2.1}{Each process has a separate memory space. A pair of \underline{base and limit registers} define the logical address space. The CPU must check every memory access generated in user mode to be sure it is between base and limit for that user.}

\separate{7}\definition{10.2.2}{The \underline{logical address space} is generated by the CPU and can also be referred to as the \underline{virtual address}.}

\separate{7}\definition{10.2.3}{The \underline{physical address space} is the address seen by the memory unit. This address cannot be seen by user applications.}

\separate{7}\definition{10.2.4}{The \underline{memory management unit} is a hardware device that maps virtual addresses to their corresponding physical ones.}

\separate{7}\definition{10.2.5}{Relocation registers are used to protect user processes from each other, and from changing operating-system code and data. This is known as \underline{memory protection}.}

\separate{7}Memory can be allocated into one of a set of holes using either: first-fit, putting the memory in the first hole that is big enough; best-fit, putting the memory in the smallest hole which is big enough; and worst-fit, putting the memory into the largest available hole.

\separate{7}\definition{10.2.6}{\underline{External fragmentation} is memory space that is able to satisfy a request but it is not contiguous.}

\separate{7}\definition{10.2.7}{\underline{Internal fragmentation} is memory that is allocated successfully, but may be slightly larger than the requested amount of memory.}

\separate{7}\definition{10.2.8}{External fragmentation can be reduced by \underline{compaction}, where memory contents are shuffled to place all free memory together in one large block.}

\separate{7}\definition{10.2.9}{Physical memory is divided into fixed-sized blocks called \underline{frame} while logical memory is divided into blocks of the same size called \underline{pages}.}

\separate{7}\definition{10.2.10}{To run a program of size \(n\) pages, we need to find \(n\) free pages before loading the program. This process is called \underline{paging}.}

\separate{7}\definition{10.2.11}{A \underline{page table} is set up to help manage the mapping of logical to physical addresses using an address translation.}

\separate{7}\definition{10.2.12}{The logical memory address generated by the CPU is divided into the \underline{page number}, which is used as an index into the page table which contains the base address of each page in physical memory, and the \underline{page offset}, which combined with the base address to define the physical memory address that is sent to the memory unit.}

\separate{7}\definition{10.2.13}{When a logical address is generated by the CPU, the MMU first checks if its page number is present in the \underline{TLB}. If the page number is found, its frame number is immediately available and is used to access memory. If the page number is not in the TLB, a memory reference to the page table has to be created. If the TLB is already full of entries, then an existing entry is selected for replacement.}

\separate{7}\definition{10.2.14}{The \underline{hit ratio} \(\alpha\), is the percentage of times that a page number is found in the associative registers.}

\separate{7}\definition{10.2.15}{The \underline{effective access time} is calculated using
\begin{equation}
    \text{EAT} = (\alpha \times \text{memory access time}) + 2((1 - \alpha) \times \text{memory access time})
\end{equation}}

\separate{7}\definition{10.2.16}{Memory protection is implemented by associating a protection bit with each friend. This is done by attaching a \underline{valid-invalid bit} to each entry where valid indicates that the associated page is in the process's logical address space, and thus is a legal page.}

\separate{7}\definition{10.2.17}{The \underline{page fault rate} is the percentage of page references which are invalid (represented as a decimal).}

\separate{7}\definition{10.2.18}{\underline{Demand paging} limits the pages in memory to those which are currently required by the process during execution. This reduces the number of pages to be processed and the number of frames allocated to a process whilst increasing overall response time for a collection of processes and the number of processes executing.}

\separate{7}\definition{10.2.19}{\underline{Page replacement} finds some page in memory that isn't really in use and swaps it out for some other required page.}

\separate{7}\definition{10.2.20}{\underline{Belady's anomaly} states that, for some page fault algorithms, the page fault rate may increase as the number of allocated frames increases.}

\separate{7}\definition{10.2.21}{\underline{Thrashing} occurs when a process is spending more time paging than doing actual work. If a process does not have enough pages, then the page-fault rate is very high.}

\separate{7}\definition{10.2.22}{\underline{Prepaging} is the process of moving all the required pages into memory at once.}

\subsection{Storage Management}

\subsection{Deadlocks}

\newpage\part{References}

\section{Module Professors}

The below is accurate as of the 2025/26 academic year and will not be updated in subsequent years.
\begin{itemize}
    \item \hyperref[part:1.de]{\textit{Digital Electronics}} and \hyperref[part:2.ma]{\textit{Machine Architecture}} are taught by Doctor Ioannis Ivrissimtzis.
    \begin{itemize}
        \item \textbf{Office Location}: Room 2023, MCS Building
        \item \textbf{Preferred Contact Methods}: ioannis.ivrissimtzis@durham.ac.uk
        \item \textbf{Pronouns}: he, him
    \end{itemize}
    
    \item \hyperref[part:3.db]{\textit{Databases}} is taught by Doctor Peter Davies-Peck.
    \begin{itemize}
        \item \textbf{Office Location}: Room 2063, MCS Building
        \item \textbf{Preferred Contact Methods}: peter.w.davies@durham.ac.uk
        \item \textbf{Pronouns}: he, him
    \end{itemize}
    
    \item \hyperref[part:4.os]{\textit{Operating Systems}} is taught by Doctor Anish Jindal.
    \begin{itemize}
        \item \textbf{Office Location}: Room 2065, MCS Building
        \item \textbf{Preferred Contact Methods}: anish.jindal@durham.ac.uk
        \item \textbf{Pronouns}: he, him
    \end{itemize}
\end{itemize}

\section{Cited Works}

\end{document}

%%%